\documentclass[12pt]{article}
\usepackage[utf8]{inputenc}
\usepackage[T1]{fontenc}
\usepackage[margin=1in]{geometry}
\usepackage{setspace}
\usepackage{xcolor}
\usepackage[colorlinks=true, urlcolor=blue!60!black, linkcolor=blue!60!black]{hyperref}
\setlength{\parindent}{0em}
\setlength{\parskip}{0.45em}
\renewcommand{\baselinestretch}{1.0}
\begin{document}
\begin{center}
{\Large \textit{Meyer v. Nebraska}, 262 U.S. 390 (1923)\par}
\vspace{0.3em}
{\normalsize Case Brief --- One-Page Edition\par}
\end{center}

\textbf{Memory Jogger.} Nebraska's ban on teaching young children modern foreign languages violated Fourteenth Amendment ``liberty'' (\href{https://www.courtlistener.com/opinion/100233/meyer-v-nebraska/}{CourtListener}).

\textbf{Facts.} Nebraska's post-World War I Siman Act barred teaching modern foreign languages to children who had not passed the eighth grade. Robert T. Meyer, a Lutheran parochial-school teacher, was charged for teaching a ten-year-old to read German from Biblical stories.

\textbf{Procedural History.} Meyer was fined \$25; the Nebraska Supreme Court affirmed, upholding the act as a police-power measure to make English the ``mother tongue'' of all children.

\textbf{Votes.} Reversed 7--2. McReynolds wrote for the Court; Holmes, joined by Sutherland, dissented in \textit{Bartels v. Iowa}.

\textbf{Holding.} The act arbitrarily infringed due process liberty and exceeded the police power, resting on the teacher's right to teach and parents' right to engage him.

\textbf{Reasoning.} Liberty reaches knowledge, marriage, child-rearing, and worship; ``[m]ere knowledge of the German language cannot reasonably be regarded as harmful.'' A wish ``to foster a homogeneous people with American ideals'' cannot justify forbidden means, and Plato's and Sparta's state-reared children are alien to American institutions. Because only modern languages were banned and no emergency existed in peacetime, the law was arbitrary. The State may still require English instruction and set public curricula. Holmes would have deferred on a question ``upon which men reasonably might differ.''

\textbf{Examples.} \emph{Same-Side:} compulsory public schooling (\textit{Pierce}). \emph{Opposite-Side:} requiring all schools to teach English. \emph{Fence-Sitter:} homeschool testing rules.

\textbf{Critique.} Originalists call it a \textit{Lochner}-era invention that seeded \textit{Griswold} and \textit{Roe}; civil libertarians praise its defense of pluralism; Barbara Bennett Woodhouse says it treated ``the child as property.''

\textbf{Key Quotation.} ``Without doubt, it denotes not merely freedom from bodily restraint, but also the right of the individual to contract, to engage in any of the common occupations of life, to acquire useful knowledge, to marry, establish a home and bring up children, to worship God according to the dictates of his own conscience, and generally to enjoy those privileges long recognized at common law as essential to the orderly pursuit of happiness by free men.'' --- Justice McReynolds, 262 U.S. at 399.

\end{document}
